\appendix

\section{Arms and blocks}\label{sec:app-arms}

\Cref{tab:arms} lists every arm with the number of its runs and what was changed against the arm it was built from.
\Cref{tab:blocks} lists every block. The same two tables are in the released files \texttt{arms.csv} and
\texttt{experiments.csv}.

\begin{table}[p]
\centering
\footnotesize
\caption{Arms of the experiments. Labels are neutral; the category names what was changed against the arm it was built from. Prompt text and bundle files are not released.}
\label{tab:arms}
\begin{tabular}{llrp{0.58\linewidth}}
\toprule
Arm & Model & Runs & What it is \\
\midrule
\texttt{BF5} & 31B & 142 & Single agent: the public Budget-Fit V5 configuration, run unchanged (\emph{reference}) \\
\texttt{BF5-p} & 31B & 30 & BF5 with one changed section of the task prompt (\emph{prompt}) \\
\texttt{BF5-c} & 31B & 32 & BF5 with a prompt-only checking step before the patch is submitted (\emph{prompt}) \\
\texttt{BF5-s} & 31B & 10 & BF5 with changed sampling settings (\emph{sampling}) \\
\texttt{W1} & 31B & 30 & Wrapped: a root agent calls a coder sub-agent as a tool; the coder has its own context (\emph{architecture}) \\
\texttt{W2} & 31B & 40 & W1 with a bounded restart of the coder after a malformed tool call (\emph{architecture}) \\
\texttt{W3} & 31B & 70 & W2 with a root that can only pass the request on and restart the coder (\emph{architecture}) \\
\texttt{W1-b} & 31B & 30 & W1 with a shorter time limit, more tool calls and half the output limit per call (\emph{budget}) \\
\texttt{W1-p} & 31B & 30 & W1 with a rewritten coder prompt, a shorter time limit and half the output limit per call (\emph{prompt}) \\
\texttt{SH1} & 31B & 30 & Wrapped, coder with a shell tool only; first version (\emph{architecture}) \\
\texttt{SH3} & 31B & 40 & SH1, third version of prompt and limits (\emph{architecture}) \\
\texttt{SH5} & 31B & 112 & SH1, fifth version of prompt and limits (\emph{architecture}) \\
\texttt{SH6} & 31B & 94 & SH5 with an added analysis step in the prompt (\emph{prompt}) \\
\texttt{ID24} & 31B & 10 & Public bundle of Team Muhammed Idris, version 24, run unchanged (\emph{reference}) \\
\texttt{ID24-r} & 31B & 10 & ID24 with one file-writing tool removed, a shorter command timeout, added prompt sections (\emph{tools}) \\
\texttt{e-single} & E4B & 24 & Single agent (\emph{architecture}) \\
\texttt{e-wrap} & E4B & 24 & Wrapped, as W1 (\emph{architecture}) \\
\texttt{e-wrap-ad} & E4B & 24 & e-wrap with a fine-tuned adapter on the coder (withheld) (\emph{adapter}) \\
\texttt{e-BF} & E4B & 64 & The public single-agent prompt adapted to E4B (six tools, 28 calls) (\emph{prompt}) \\
\texttt{e-BF-s} & E4B & 16 & e-BF with changed sampling settings (\emph{sampling}) \\
\texttt{base} & E4B & 1{,}848 & The withheld E4B agent bundle; the baseline of every later comparison (\emph{baseline}) \\
\texttt{base-p} & E4B & 16 & base with one prompt change (\emph{prompt}) \\
\texttt{v-p1} & E4B & 96 & base with one prompt-only change (\emph{prompt}) \\
\texttt{v-p2} & E4B & 384 & base with one prompt-only change (candidate 1) (\emph{prompt}) \\
\texttt{v-p3} & E4B & 96 & base with one prompt-only change (\emph{prompt}) \\
\texttt{v-p4} & E4B & 96 & base with one prompt-only change (\emph{prompt}) \\
\texttt{v-b1} & E4B & 96 & base with a changed thinking-token limit (\emph{budget}) \\
\texttt{v-t1} & E4B & 15 & base with a changed file-editing path, first build (\emph{tools}) \\
\texttt{v-t2} & E4B & 15 & as v-t1, second design, first build (\emph{tools}) \\
\texttt{v-t3} & E4B & 15 & rebuild of v-t1, option a (\emph{tools}) \\
\texttt{v-t4} & E4B & 111 & rebuild of v-t1, option b (\emph{tools}) \\
\texttt{v-t5} & E4B & 15 & rebuild of v-t2, option a (\emph{tools}) \\
\texttt{v-t6} & E4B & 623 & rebuild of v-t2, option b (candidate 2) (\emph{tools}) \\
\texttt{v-a1} & E4B & 16 & base with an added sub-agent step, first design, first build (\emph{architecture}) \\
\texttt{v-a2} & E4B & 16 & as v-a1, second design, first build (\emph{architecture}) \\
\texttt{v-a3} & E4B & 608 & rebuild of v-a1 (candidate 3) (\emph{architecture}) \\
\texttt{v-a4} & E4B & 96 & rebuild of v-a2 (\emph{architecture}) \\
\bottomrule
\end{tabular}
\end{table}


\begin{table}[p]
\centering
\footnotesize\setlength{\tabcolsep}{4.5pt}
\caption{Every block of agent runs in the released file \texttt{runs.csv}. A block is one session on one machine; its arms ran in one job queue. Machine: one GPU with emulated L4 latency (emulated), a hosted notebook with four L4 GPUs (4$\times$L4), or the single-GPU workstation. Source: the development tasks of the competition (dev.) or Berth. Rep.: repeats per task and arm.}
\label{tab:blocks}
\begin{tabular}{llllrrrrp{0.27\linewidth}}
\toprule
Block & Model & Machine & Source & Tasks & Arms & Rep. & Runs & Note \\
\midrule
\texttt{pub31-1} & 31B & emulated & dev. & 30 & 3 & 1 & 90 & first 30 development tasks \\
\texttt{pub31-2} & 31B & emulated & dev. & 30 & 1 & 1 & 30 & same 30 tasks \\
\texttt{pub31-3} & 31B & emulated & dev. & 30 & 1 & 1 & 30 & same 30 tasks \\
\texttt{pub31-4} & 31B & emulated & dev. & 30 & 2 & 1 & 60 & same 30 tasks, both arms in one queue \\
\texttt{pub31-5} & 31B & emulated & dev. & 40 & 2 & 1 & 80 & 40 further development tasks \\
\texttt{pub31-6} & 31B & emulated & dev. & 40 & 2 & 1 & 80 & same 40 tasks, both arms in one queue \\
\texttt{pub31-7} & 31B & emulated & dev. & 70 & 1 & 1 & 70 & all 70 tasks \\
\texttt{pub31-8} & 31B & emulated & dev. & 70 & 3 & 1 & 130 & all 70 tasks; 30 tasks with masked test runners \\
\texttt{l4-1} & 31B & 4$\times$L4 & dev. & 10 & 1 & 1 & 12 & 10 tasks, plus 2 masked \\
\texttt{l4-2} & 31B & 4$\times$L4 & dev. & 24 & 1 & 1 & 24 & 24 tasks \\
\texttt{pubE4-1} & E4B & 4$\times$L4 & dev. & 24 & 3 & 1 & 72 & 24 development tasks with small fixes \\
\texttt{berth31-time} & 31B & 4$\times$L4 & Berth & 10 & 4 & 1 & 40 & 10 tasks, serial \\
\texttt{berth31-check} & 31B & 4$\times$L4 & Berth & 16 & 2 & 2 & 64 & 16 tasks \\
\texttt{berthE4-k1} & E4B & 4$\times$L4 & Berth & 24 & 2 & 2 & 96 & 24 tasks \\
\texttt{berthE4-k2} & E4B & 4$\times$L4 & Berth & 8 & 4 & 2 & 64 & 8 tasks \\
\texttt{parity-1} & E4B & workstation & Berth & 24 & 1 & 1 & 24 & no sandbox, 8 workers \\
\texttt{parity-2} & E4B & workstation & Berth & 24 & 1 & 1 & 24 & no sandbox, 1 worker \\
\texttt{parity-3} & E4B & workstation & Berth & 24 & 1 & 1 & 24 & container runner, version 1 \\
\texttt{parity-4} & E4B & workstation & Berth & 24 & 1 & 1 & 24 & container runner, version 3 \\
\texttt{parity-5} & E4B & workstation & Berth & 24 & 1 & 1 & 24 & container runner, version 4, run a \\
\texttt{parity-6} & E4B & workstation & Berth & 24 & 1 & 1 & 24 & container runner, version 4, run b \\
\texttt{parity-7} & E4B & workstation & Berth & 24 & 1 & 1 & 24 & container runner, version 7 (final) \\
\texttt{screen-1} & E4B & workstation & Berth & 48 & 6 & 2 & 576 & screen set S \\
\texttt{confirm-1} & E4B & workstation & Berth & 72 & 2 & 4 & 576 & small panel C \\
\texttt{pilot-2a} & E4B & workstation & Berth & 15 & 2 & 1 & 30 & 15 tasks of the case-study panel \\
\texttt{pilot-2b} & E4B & workstation & Berth & 15 & 4 & 1 & 60 & same 15 tasks \\
\texttt{screen-2} & E4B & workstation & Berth & 48 & 3 & 2 & 288 & screen set S \\
\texttt{confirm-2} & E4B & workstation & Berth & 128 & 2 & 4 & 1{,}024 & panel H \\
\texttt{dev-3} & E4B & workstation & Berth & 16 & 3 & 1 & 48 & 16 tasks of panel C \\
\texttt{screen-3} & E4B & workstation & Berth & 48 & 3 & 2 & 288 & screen set S \\
\texttt{confirm-3} & E4B & workstation & Berth & 128 & 2 & 4 & 1{,}024 & panel H, second use \\
\bottomrule
\end{tabular}
\end{table}


\section{Recomputation of the case study with the released kit}\label{sec:app-recompute}

% written by scripts/reuse_short.py from the sources of the short paper; do not edit
The released table \texttt{e8\_compare\_table.csv} is in the input format of the \texttt{compare} module of the kit.
From a clone of the repository, two commands recompute the analysis:

\begin{verbatim}
pip install "./drydock-eval[compare]"
drydock-eval compare signflip e8_compare_table.csv --reference F --repeats 4
\end{verbatim}

We ran this command from the released source tree with the script \texttt{supp/reproduce.py}, which stops on any
mismatch. For all four arms the command returned the registered estimate, the exact $p$ as the same fraction, the
Holm-adjusted $p$ and every repository-level quantity of Table~\ref{tab:s-repolevel}.

Two quantities depend on Monte Carlo draws: the inverted intervals, and the Monte Carlo $p$ that is reported next to
the exact one. The draws are frozen (seed \SReproSeed), but which draw meets which task depends on the order of the
tasks. The registered analysis used the order of the panel file; the command sorts the task identifiers. With the
tasks in panel order, the routine of the kit returns the registered values to the last digit
(Table~\ref{tab:s-repro}). With the order of the command, the ends of the 95\% intervals move by at most
\SReproMoveNinetyFive{} points and those of the simultaneous intervals by at most \SReproMoveSim{} points. The largest
move is the upper end of the interval of H, $+\SReproHCmdHi$ instead of $+\SReproHRegHi$ points. The $p$-value of the
test is a step function of the shift, and across that stretch it stays between \SReproStepLo{} and \SReproStepHi{},
next to the 0.05 it is compared with. Both intervals include zero. The paper reports the registered intervals.

\begin{table}[!htb]
\centering
\caption{Reproduction with the released kit: the two quantities that depend on Monte Carlo draws. 95\% intervals of the change against F in percentage points, and the Monte Carlo $p$ that accompanies the exact $p$.}\label{tab:s-repro}
\footnotesize
\begin{tabular}{lrrrrrr}
\toprule
 & \multicolumn{3}{c}{95\% interval} & \multicolumn{3}{c}{Monte Carlo $p$} \\
\cmidrule(lr){2-4}\cmidrule(lr){5-7}
Arm & Registered & Kit, panel order & Kit command & Registered & Kit, panel order & Kit command \\
\midrule
R & [${-}2.82$, ${+}2.17$] & [${-}2.82$, ${+}2.17$] & [${-}2.80$, ${+}2.16$] & 0.8650 & 0.8650 & 0.8642 \\
C & [${-}2.03$, ${+}2.48$] & [${-}2.03$, ${+}2.48$] & [${-}2.04$, ${+}2.48$] & 0.9243 & 0.9243 & 0.9268 \\
H & [${-}4.45$, ${+}0.38$] & [${-}4.45$, ${+}0.38$] & [${-}4.45$, ${+}0.25$] & 0.1106 & 0.1106 & 0.1095 \\
T & [${-}1.42$, ${+}2.94$] & [${-}1.42$, ${+}2.94$] & [${-}1.42$, ${+}2.94$] & 0.5599 & 0.5599 & 0.5624 \\
\bottomrule
\end{tabular}
\end{table}



\section{The three decisions recomputed from the released rows}\label{sec:app-recheck}

\begin{table}[t]
\centering
\footnotesize\setlength{\tabcolsep}{3.6pt}
\caption{The three applications of the pass rule (four repeats per task and arm). Gain: candidate minus baseline in percentage points. $p$: one-sided, repository-clustered. $+/-$: repositories with a positive or negative summed gain. LOO: lowest gain when one repository is left out. Time: ratio of mean agent time, candidate to baseline.}
\label{tab:confirm}
\begin{tabular}{rlrrrrrrrrrrl}
\toprule
No. & Candidate & Tasks & Repos & Runs/arm & Base & Cand. & Gain & $p$ & $+/-$ & LOO & Time & Verdict \\
\midrule
1 & \texttt{v-p2} & 72 & 18 & 288 & 63 & 62 & \ensuremath{-}0.35 & 0.6034 & 3/6 & \ensuremath{-}1.47 & 1.03 & not confirmed \\
2 & \texttt{v-t6} & 128 & 19 & 512 & 103 & 101 & \ensuremath{-}0.39 & 0.6034 & 7/6 & \ensuremath{-}0.85 & 1.18 & not confirmed \\
3 & \texttt{v-a3} & 128 & 19 & 512 & 97 & 118 & +4.10 & 0.0405 & 9/6 & +2.94 & 1.10 & not confirmed \\
\bottomrule
\end{tabular}
\end{table}


The frozen reports of the three applications of the pass rule were written by the registered analysis scripts on the
machine that ran the jobs, from result files that are not released. The released file \texttt{runs.csv} holds the
same runs with a reduced set of columns. The script \texttt{scripts/recheck\_gates.py} builds, for each application,
the input table of the \texttt{compare} module from these rows, calls the routines of the kit at tag
\texttt{v0.4.0} with the registered thresholds and seed, and compares the result with the released copy of the
frozen report, \texttt{pass\_rule\_applications.json}. It stops on the first mismatch.

All \RecheckQuantities{} compared quantities agree: for each application the solved counts of both arms, the gain,
the bootstrap $p$ with its 9,999 draws, the gain and the leave-one-out estimate of every repository, the counts of
repositories that gain and lose, and the decision. \Cref{tab:confirm} gives the values. The time ratio is the one
quantity that the recomputation can only approximate, because the reports computed it from timings with more
digits than the released rows keep; recomputed from the released rows it is \RecheckTimeOne{}, \RecheckTimeTwo{} and
\RecheckTimeThree{}, which round to the reported values.

\section{Released data}\label{sec:app-data}

\begin{table}[t]
\centering
\footnotesize
\caption{Data files released with this paper, on the branch \texttt{paper} of the repository, in the directory
\texttt{data/}.}
\label{tab:datafiles}
\begin{tabular}{@{}lrp{0.56\linewidth}@{}}
\toprule
File & Records & Content \\
\midrule
\texttt{runs.csv} & \RunsTotal & one row per agent run of \cref{sec:big,sec:small,sec:programmes} \\
\texttt{controls.csv} & \ControlRows & one row per control grading (reference patch or empty patch) \\
\texttt{serving\_bench.csv} & \ServeRows & one row per cell of the serving benchmark of \cref{sec:serving} \\
\texttt{serving\_servers.csv} & 7 & start-up time and cache size of each model server of the benchmark \\
\texttt{experiments.csv} & \BlocksTotal & one row per block (\cref{tab:blocks}) \\
\texttt{arms.csv} & \ArmsTotal & one row per arm (\cref{tab:arms}) \\
\texttt{submissions.csv} & 5 & our hosted submissions with their displayed scores \\
\texttt{pass\_rule\_applications.json} & 3 & the frozen report of each application of the pass rule \\
\texttt{e8\_runs.csv} & \CaseRows & one row per job of the case study, with the rows of interrupted jobs \\
\texttt{e8\_tasks.csv} & 231 & the tasks of the case study with their solved counts per arm \\
\texttt{e8\_repositories.csv} & 15 & the same per repository \\
\texttt{e8\_compare\_table.csv} & 231 & the case study in the input format of the \texttt{compare} module \\
\texttt{e8\_summary.json} & 1 & the estimates, intervals and $p$-values of the case study \\
\bottomrule
\end{tabular}
\end{table}

\Cref{tab:datafiles} lists the files. \texttt{data/README.md} describes every column, and \texttt{data/SHA256SUMS}
holds the checksums. The file \texttt{runs.csv} has 29 columns:
\begin{description}[leftmargin=1.2em,itemsep=0pt,topsep=2pt,font=\normalfont\ttfamily]
\item[experiment, block] the experiment family and the block (\cref{tab:blocks}).
\item[model, machine, task\_source] 31B or E4B; the machine of \cref{sec:setting}; development tasks or Berth.
\item[arm, condition] the arm (\cref{tab:arms}); \texttt{standard}, or \texttt{masked\_test\_runners} for the runs
of \cref{tab:masked}.
\item[task, repository, repeat] the task (an opaque code for development tasks, the Berth identifier otherwise), its
repository, and the number of the repeat within the block.
\item[resolved] 1 if the tests of the task passed on the patch that the run left behind, else 0.
\item[end\_reason, outcome\_id] how the run ended (normally, at the time limit, at the tool-call limit, without a
submission), and the outcome class of the runner.
\item[wall\_s, t\_agent\_s, t\_setup\_s, t\_verify\_s] seconds of the whole job and of its three phases, where the
harness recorded them.
\item[emulation\_pad\_s] seconds of padding added to model calls on the emulated-latency machine.
\item[model\_calls, model\_calls\_failed, tool\_calls] counts per run.
\item[tokens\_in, tokens\_out, max\_prompt\_tokens] prompt tokens summed over the requests of the run, generated
tokens, and the longest prompt.
\item[summaries] number of history summaries that the harness made in the run.
\item[patch\_chars, test\_exit] length of the final patch, and the exit status of the test command.
\item[start\_utc, end\_utc] times of the job, where recorded.
\end{description}
No column holds prompt text, model output, error text, a path or a machine name. The columns of
\texttt{e8\_runs.csv} are listed in \cref{tab:s-columns}.

\begin{table}[!htb]
\centering
\caption{Columns of \texttt{e8\_runs.csv}. The rows of interrupted jobs hold the first ten columns and the two times of day only.}\label{tab:s-columns}
\footnotesize
\begin{tabular}{lp{11.4cm}}
\toprule
Column & Content \\
\midrule
\texttt{batch} & 1 to 5: the batch of the block in which the job ran; batch 5 holds the re-runs \\
\texttt{stage} & 1, 2 or 3: the nested stage of the job list that holds the task (118, 37 and 76 tasks) \\
\texttt{arm} & F, R, C, H or T \\
\texttt{task} & task identifier in the public task collection the panel was drawn from \\
\texttt{repository} & repository of the task (owner/name); the unit of clustering \\
\texttt{repeat} & 1 to 4 \\
\texttt{attempt} & 1; 2 for the re-run of an interrupted job \\
\texttt{row\_class} & agent: an analysed run. infrastructure: an interrupted job, not analysed \\
\texttt{outcome\_id} & A00 solved; A01 unsolved, run completed; A02 tool-call budget reached; A05 unsolved at the five-minute limit; A06 no submission; A18 unsolved, over-long prompt rejected; I09 interrupted job \\
\texttt{resolved} & 1 if the run is solved (the tests of the task pass on the edits the run left), else 0 \\
\texttt{t\_agent\_s} & seconds of the agent phase of the job \\
\texttt{t\_setup\_s} & seconds of the setup phase, before the agent starts \\
\texttt{t\_verify\_s} & seconds of the grading phase; empty if no test was run \\
\texttt{model\_calls} & chat requests of the run that returned a response \\
\texttt{model\_calls\_failed} & chat requests of the run that raised an error instead \\
\texttt{tool\_calls} & tool calls made by the agent \\
\texttt{tokens\_in} & prompt tokens, summed over the requests of the run \\
\texttt{tokens\_out} & completion tokens, summed over the requests of the run \\
\texttt{max\_prompt\_tokens} & largest prompt of the run, in tokens \\
\texttt{summaries} & history summaries (compactions) made during the run \\
\texttt{patch\_chars} & characters of the diff the run left; 0 if none \\
\texttt{test\_exit} & exit status of the test command; 0 = tests passed; $-1$ = no test was run \\
\texttt{start\_utc} & start of the job (UTC) \\
\texttt{end\_utc} & end of the job (UTC) \\
\texttt{requests\_checked} & returned requests of the run that the frozen checker re-rendered \\
\texttt{requests\_changed} & of them, requests whose prompt the reference rule of arm F would render differently (the exposure count; 0 in F) \\
\texttt{check\_passed} & 1 if the run passed the prompt check \\
\bottomrule
\end{tabular}
\end{table}

